\section{PBRS e imitación}
Se combinan recompensas potenciales y aprendizaje por imitación para orientar
la exploración y aprovechar demostraciones antes del ajuste con RL.

\subsection{Shaping potencial e invariancia}
PBRS \cite{ng1999} añade a la recompensa base una diferencia de potenciales:
\[
    r_t^{\rm sh}=r_t+F_t,\qquad F_t=\gamma\Phi(s_{t+1})-\Phi(s_t).
\]
El entorno permite escalar este término por una constante, que puede absorberse
en el potencial. Se usan los potenciales de distancia al balón, avance controlado
y proximidad al arco definidos en el modelado del entorno. En toda terminación,
incluido el horizonte, el potencial siguiente se fija en cero.

Para un episodio de longitud $T$, la suma descontada telescopa:
\begin{equation}
    \sum_{t=0}^{T-1}\gamma^t F_t=-\Phi(s_0).
\end{equation}
Por tanto, el retorno y el valor óptimo solo cambian en una constante dependiente
del estado inicial:
\[
    G^{\rm sh}=G-\Phi(s_0).
\]
\[
    Q_{\rm sh}^*(s,a)=Q^*(s,a)-\Phi(s).
\]
Como el desplazamiento no depende de la acción, se conserva la política óptima.
Esta garantía corresponde al problema de horizonte finito y al descuento usado
en la suma; no garantiza por sí sola la convergencia de una red neuronal.

La alternativa ingenua recompensa cualquier avance horizontal positivo:
\[
    F_t^{\rm ingenuo}=\max(0,x_{b,t+1}-x_{b,t}).
\]
Puede premiar avances repetidos sin compensar los retrocesos ni exigir control
del balón. El contraste usa el mismo modelo inicial de conducción y cambia
únicamente el shaping entre ninguno, PBRS e ingenuo para estudiar reward hacking.

\subsection{Preentrenamiento por imitación}
Las demostraciones asocian observaciones con acciones expertas, humanas o
programadas. Behavioral Cloning (BC) inicializa las mismas redes que luego se
ajustan con RL. Para PPO/A2C, la política minimiza la entropía cruzada:
\[
    L_{\rm BC}=-\mathbb{E}_{(s,a^E)\sim\mathcal{D}}[\log\pi_\theta(a^E\mid s)].
\]
El crítico se ajusta a los retornos descontados de las trayectorias registradas.
La implementación permite balancear clases y añadir regularización entrópica.

Para DQN se utiliza un margen supervisado inspirado en DQfD \cite{hester2018}:
\[
    L_{\rm margen}=\max_a[Q_\theta(s,a)+\ell(a^E,a)]-Q_\theta(s,a^E).
\]
El margen es cero para la acción experta y positivo para las restantes. Se añade
una pérdida TD de Double DQN, ponderada y calculada sobre transiciones válidas;
la supervisión de acciones utiliza únicamente estados etiquetados.

\subsection{DAgger y ajuste con RL}
Con el experto programado, DAgger \cite{ross2011} deja actuar a la red y etiqueta
cada estado visitado. Al finalizar cada ronda, incorpora las nuevas etiquetas
y reentrena sobre el conjunto acumulado:
\[
    \mathcal{D}_{k+1}=\mathcal{D}_k\cup\{(s,\pi_E(s)):s\text{ visitado}\}.
\]
Así se corrigen errores en estados inducidos por la propia política. En la variante
humana HG-DAgger \cite{kelly2019}, el usuario toma el control para corregirla;
solo esas intervenciones aportan etiquetas de acción.

El experto de conducción persigue, disputa, esquiva o pasa según la posición
del defensor y la disponibilidad del compañero. El del coordinador elige
persecución si el balón está a más de cinco metros; en otro caso selecciona
definición al cumplirse la zona de traspaso, incluida la ventaja opcional,
y conducción en los demás estados.

Finalmente se continúa el aprendizaje con RL desde los pesos imitados. En la
configuración base del sistema jerárquico, PPO utiliza tasa de aprendizaje
$10^{-4}$, coeficiente de entropía $10^{-3}$ y KL objetivo 0,02; otras variantes
ajustan estos parámetros. Los resultados de esta aceleración se evalúan después.
